% Partie III condensée (FR) — état de l'art NARRATIF. Il raconte l'histoire du champ
% qui pose la situation initiale de la thèse ; les citations (réelles) servent le
% récit, elles ne s'empilent pas.

\noindent Le traitement automatique des langues a convergé, ces dernières années, vers une même
recette : pré-entraîner un modèle sur de grandes quantités de texte, puis l'adapter à
une tâche. Les modèles de type BERT~\citep{devlin_bert_2019}, fondés sur le
Transformer~\citep{vaswani2017attention}, ont imposé cette recette, et
CamemBERT~\citep{martin-etal-2020-camembert} l'a portée au français ; la montée en
échelle des modèles génératifs~\citep{gpt3,touvron2023llama} en est le prolongement.
Tout cela repose sur une même condition : du texte, en abondance, dans le domaine visé.
Spécialiser un modèle à la médecine se fait alors naturellement par un pré-entraînement
continu sur du texte biomédical~\citep{gururangan_dont_2020}, comme l'ont montré
BioBERT~\citep{lee_biobert_2019} puis PubMedBERT~\citep{gu_domain-specific_2022}.

C'est précisément cette condition qui manque au domaine clinique. Le texte dont le modèle
aurait le plus besoin, le compte-rendu hospitalier, est aussi celui qu'il ne peut pas
obtenir : donnée personnelle, réglementée en France par la CNIL, il ne quitte pas
l'hôpital, et un modèle entraîné en son sein ne peut en sortir non plus. Les
adaptations françaises qui ont réussi l'ont fait sur des dossiers hospitaliers privés,
et sont restées privées~\citep{dura_learning_2022} ; les tentatives publiques se sont
appuyées sur des corpus restreints~\citep{copara_contextualized_2020,le_clercq_de_lannoy_strategies_2022}.
Le champ se trouve alors devant une alternative : renoncer à la cible clinique, ou tout
reconstruire à partir de texte public.

Mais le texte biomédical public n'est pas prêt à l'emploi. Il est dispersé et
massivement anglophone~\citep{labrak-etal-2024-biomistral}, et la machinerie qui a rendu
possible la curation à l'échelle du web~\citep{penedo2024fineweb}, comme les méthodes
récentes de génération de texte d'entraînement synthétique~\citep{yang2025entigraph},
ont été conçues pour l'anglais et le domaine général. La curation biomédicale existante
filtre au niveau du document~\citep{chen2023meditron}, sans distinguer le pertinent du
générique à l'intérieur d'un article. Et la connaissance structurée dont dépendent
les tâches cliniques reste sous-exploitée : la recherche est dominée par UMLS~\citep{lindberg_unified_1993},
alors que les hôpitaux français codent avec des ontologies nationales, la CIM-10
maintenue par l'ATIH~\citep{atih_pmsi}, que nous n'avons vu aucun pré-entraînement exploiter.

La tâche visée, enfin, résiste par sa structure même. L'essentiel de l'information
clinique est dans les comptes-rendus~\citep{raghavan_how_2014}, mais le codage médical
compte des dizaines de milliers d'étiquettes : un classifieur à tête fixe ne peut
couvrir des codes jamais vus à l'entraînement. Les extracteurs à vocabulaire ouvert,
GLiNER~\citep{zaratiana2024gliner} et UniversalNER~\citep{zhou2024universalner}, lèvent
cette limite en décrivant chaque étiquette en texte libre, mais, entraînés sur du texte
web, ils perdent en précision sur l'écrit clinique et les codes rares, y compris leur
variante biomédicale~\citep{zaratiana2025glinerbiomed}. L'évaluation elle-même n'est pas
stabilisée~\citep{drbenchmark2024}, les scores dépendant du tokeniseur. Pourtant, pour
ces tâches d'extraction, les encodeurs restent compétitifs et assez petits pour tourner
là où les données doivent rester~\citep{lehman_we_2023}.

Telle est la situation dont part cette thèse. Construire un extracteur clinique
partageable, sans jamais voir de dossier réel, oblige à reconstruire toute la chaîne
depuis le texte public : le corpus, le modèle, la supervision et, à la fin,
l'évaluation elle-même.
